\documentclass[10pt]{article}
% Preambulo compartido por los entregables escritos de Diablos AeroVision (HACK-018).
% Compila con pdflatex (TeX Live, MiKTeX u Overleaf). Ver docs/latex/README.md.
\usepackage[T1]{fontenc}
\usepackage[utf8]{inputenc}
\usepackage[spanish,es-noshorthands,es-tabla,es-nodecimaldot]{babel}
\usepackage{lmodern}
\usepackage{microtype}
\usepackage[a4paper,margin=2.2cm]{geometry}
\usepackage{amsmath,amssymb}
\usepackage{booktabs}
\usepackage{tabularx}
\usepackage{graphicx}
\usepackage{float}
\usepackage{enumitem}
\usepackage[table]{xcolor}
\usepackage{caption}
\usepackage{fancyhdr}
\usepackage[hidelinks]{hyperref}

\definecolor{diablos}{HTML}{B3121B}
\definecolor{grisfondo}{HTML}{F4F4F5}
\captionsetup{font=small,labelfont=bf}
\setlist{itemsep=2pt,topsep=4pt}
\setlength{\parskip}{4pt}

% \cf{clave}: cifra leida de reports/ via generated/cifras.tex. Una clave inexistente
% se imprime en rojo para que nunca pase inadvertida (y el test la rechaza).
\makeatletter
\newcommand{\cf}[1]{\@ifundefined{cf@#1}{\textcolor{red}{\textbf{??#1}}}{\@nameuse{cf@#1}}}
\makeatother
% ARCHIVO GENERADO por analysis/build_paper_numbers.py. No editar a mano.
\makeatletter
\@namedef{cf@ca.changeup.sc}{91.5}
\@namedef{cf@ca.changeup.sd}{\ensuremath{-}0.51}
\@namedef{cf@ca.changeup.sd.abs}{0.51}
\@namedef{cf@ca.changeup.sn}{92.0}
\@namedef{cf@ca.changeup.uso}{28}
\@namedef{cf@ca.slider.sc}{96.8}
\@namedef{cf@ca.slider.sd}{+0.20}
\@namedef{cf@ca.slider.sd.abs}{0.20}
\@namedef{cf@ca.slider.sn}{96.6}
\@namedef{cf@ca.slider.uso}{72}
\@namedef{cf@ca1.ct}{+0.85}
\@namedef{cf@ca1.ct.abs}{0.85}
\@namedef{cf@ca1.gb}{59}
\@namedef{cf@ca1.id}{pitcher\_00778}
\@namedef{cf@ca1.ln}{102.1}
\@namedef{cf@ca1.mano}{zurdo}
\@namedef{cf@ca1.n}{1,177}
\@namedef{cf@ca1.pc}{100.6}
\@namedef{cf@ca1.pn}{101.5}
\@namedef{cf@ca1.rank.delta}{5}
\@namedef{cf@ca1.rank.sn}{423}
\@namedef{cf@ca1.rp}{+1.02}
\@namedef{cf@ca1.rp.abs}{1.02}
\@namedef{cf@ca1.sc}{96.1}
\@namedef{cf@ca1.sd}{+0.31}
\@namedef{cf@ca1.sd.abs}{0.31}
\@namedef{cf@ca1.slot}{33}
\@namedef{cf@ca1.sn}{95.8}
\@namedef{cf@ca1.velo}{82.7}
\@namedef{cf@ca1.year}{2024}
\@namedef{cf@ca2.ct}{+1.45}
\@namedef{cf@ca2.ct.abs}{1.45}
\@namedef{cf@ca2.gb}{74}
\@namedef{cf@ca2.id}{pitcher\_00778}
\@namedef{cf@ca2.ln}{102.3}
\@namedef{cf@ca2.mano}{zurdo}
\@namedef{cf@ca2.n}{1,280}
\@namedef{cf@ca2.pc}{100.3}
\@namedef{cf@ca2.pn}{101.3}
\@namedef{cf@ca2.rank.delta}{11}
\@namedef{cf@ca2.rank.sn}{432}
\@namedef{cf@ca2.rp}{+0.10}
\@namedef{cf@ca2.rp.abs}{0.10}
\@namedef{cf@ca2.sc}{95.3}
\@namedef{cf@ca2.sd}{+0.00}
\@namedef{cf@ca2.sd.abs}{0.00}
\@namedef{cf@ca2.slot}{34}
\@namedef{cf@ca2.sn}{95.3}
\@namedef{cf@ca2.velo}{81.8}
\@namedef{cf@ca2.year}{2025}
\@namedef{cf@cal.max.global}{0.022}
\@namedef{cf@cal.stuff.park.mean}{0.0105}
\@namedef{cf@canal.n}{778}
\@namedef{cf@canal.park}{\ensuremath{-}0.0060}
\@namedef{cf@canal.park.abs}{0.0060}
\@namedef{cf@canal.park.ci}{[\ensuremath{-}0.0079, \ensuremath{-}0.0040]}
\@namedef{cf@canal.park.per100}{\ensuremath{-}0.60}
\@namedef{cf@canal.park.per100.abs}{0.60}
\@namedef{cf@canal.park.share}{42}
\@namedef{cf@canal.ratio}{4.7}
\@namedef{cf@canal.stuff}{\ensuremath{-}0.0013}
\@namedef{cf@canal.stuff.abs}{0.0013}
\@namedef{cf@canal.stuff.ci}{[\ensuremath{-}0.0013, \ensuremath{-}0.0012]}
\@namedef{cf@canal.stuff.per100}{\ensuremath{-}0.13}
\@namedef{cf@canal.stuff.per100.abs}{0.13}
\@namedef{cf@canal.stuff.share}{9}
\@namedef{cf@canal.total}{\ensuremath{-}0.0141}
\@namedef{cf@canal.total.abs}{0.0141}
\@namedef{cf@canal.total.ci}{[\ensuremath{-}0.0169, \ensuremath{-}0.0117]}
\@namedef{cf@canal.total.per100}{\ensuremath{-}1.41}
\@namedef{cf@canal.total.per100.abs}{1.41}
\@namedef{cf@canal.total.share}{100}
\@namedef{cf@canal.unallocated}{\ensuremath{-}0.0069}
\@namedef{cf@canal.unallocated.abs}{0.0069}
\@namedef{cf@canal.unallocated.ci}{[\ensuremath{-}0.0086, \ensuremath{-}0.0052]}
\@namedef{cf@canal.unallocated.per100}{\ensuremath{-}0.69}
\@namedef{cf@canal.unallocated.per100.abs}{0.69}
\@namedef{cf@canal.unallocated.share}{48}
\@namedef{cf@cb.ct}{+2.49}
\@namedef{cf@cb.ct.abs}{2.49}
\@namedef{cf@cb.fourseam.sc}{95.3}
\@namedef{cf@cb.fourseam.sd}{\ensuremath{-}9.03}
\@namedef{cf@cb.fourseam.sd.abs}{9.03}
\@namedef{cf@cb.fourseam.sn}{104.3}
\@namedef{cf@cb.fourseam.uso}{77}
\@namedef{cf@cb.gb}{56}
\@namedef{cf@cb.id}{pitcher\_00204}
\@namedef{cf@cb.ln}{99.1}
\@namedef{cf@cb.mano}{zurdo}
\@namedef{cf@cb.n}{658}
\@namedef{cf@cb.pc}{98.5}
\@namedef{cf@cb.pctl}{97}
\@namedef{cf@cb.pn}{100.2}
\@namedef{cf@cb.rank.delta}{461}
\@namedef{cf@cb.rank.sn}{13}
\@namedef{cf@cb.rp}{+1.05}
\@namedef{cf@cb.rp.abs}{1.05}
\@namedef{cf@cb.sc}{99.5}
\@namedef{cf@cb.sd}{\ensuremath{-}6.99}
\@namedef{cf@cb.sd.abs}{6.99}
\@namedef{cf@cb.sinker.sc}{102.0}
\@namedef{cf@cb.sinker.sd}{\ensuremath{-}1.40}
\@namedef{cf@cb.sinker.sd.abs}{1.40}
\@namedef{cf@cb.sinker.sn}{103.4}
\@namedef{cf@cb.sinker.uso}{0}
\@namedef{cf@cb.slider.sc}{113.8}
\@namedef{cf@cb.slider.sd}{\ensuremath{-}0.05}
\@namedef{cf@cb.slider.sd.abs}{0.05}
\@namedef{cf@cb.slider.sn}{113.9}
\@namedef{cf@cb.slider.uso}{22}
\@namedef{cf@cb.slot}{41}
\@namedef{cf@cb.sn}{106.4}
\@namedef{cf@cb.velo}{93.5}
\@namedef{cf@cb.year}{2025}
\@namedef{cf@cc.canalparquebip}{\ensuremath{-}1.86}
\@namedef{cf@cc.canalparquebip.abs}{1.86}
\@namedef{cf@cc.canaltotal}{\ensuremath{-}4.41}
\@namedef{cf@cc.canaltotal.abs}{4.41}
\@namedef{cf@cc.changeup.sc}{105.1}
\@namedef{cf@cc.changeup.sd}{\ensuremath{-}1.68}
\@namedef{cf@cc.changeup.sd.abs}{1.68}
\@namedef{cf@cc.changeup.sn}{106.8}
\@namedef{cf@cc.changeup.uso}{16}
\@namedef{cf@cc.ct}{\ensuremath{-}4.41}
\@namedef{cf@cc.ct.abs}{4.41}
\@namedef{cf@cc.curveball.sc}{107.6}
\@namedef{cf@cc.curveball.sd}{+0.29}
\@namedef{cf@cc.curveball.sd.abs}{0.29}
\@namedef{cf@cc.curveball.sn}{107.3}
\@namedef{cf@cc.curveball.uso}{11}
\@namedef{cf@cc.fourseam.sc}{90.3}
\@namedef{cf@cc.fourseam.sd}{\ensuremath{-}9.77}
\@namedef{cf@cc.fourseam.sd.abs}{9.77}
\@namedef{cf@cc.fourseam.sn}{100.1}
\@namedef{cf@cc.fourseam.uso}{73}
\@namedef{cf@cc.gb}{46}
\@namedef{cf@cc.id}{pitcher\_00466}
\@namedef{cf@cc.ln}{99.8}
\@namedef{cf@cc.mano}{derecho}
\@namedef{cf@cc.n}{868}
\@namedef{cf@cc.pc}{98.7}
\@namedef{cf@cc.pn}{100.6}
\@namedef{cf@cc.rank.delta}{462}
\@namedef{cf@cc.rank.sn}{141}
\@namedef{cf@cc.rp}{+0.15}
\@namedef{cf@cc.rp.abs}{0.15}
\@namedef{cf@cc.sc}{94.5}
\@namedef{cf@cc.sd}{\ensuremath{-}7.40}
\@namedef{cf@cc.sd.abs}{7.40}
\@namedef{cf@cc.slot}{48}
\@namedef{cf@cc.sn}{101.9}
\@namedef{cf@cc.velo}{91.3}
\@namedef{cf@cc.year}{2025}
\@namedef{cf@cf.changeup.cdmx}{101.5}
\@namedef{cf@cf.changeup.delta}{\ensuremath{-}1.6}
\@namedef{cf@cf.changeup.delta.abs}{1.6}
\@namedef{cf@cf.changeup.neutral}{103.1}
\@namedef{cf@cf.curveball.cdmx}{108.2}
\@namedef{cf@cf.curveball.delta}{\ensuremath{-}0.5}
\@namedef{cf@cf.curveball.delta.abs}{0.5}
\@namedef{cf@cf.curveball.neutral}{108.7}
\@namedef{cf@cf.cutter.cdmx}{95.2}
\@namedef{cf@cf.cutter.delta}{\ensuremath{-}0.4}
\@namedef{cf@cf.cutter.delta.abs}{0.4}
\@namedef{cf@cf.cutter.neutral}{95.6}
\@namedef{cf@cf.fourseam.cdmx}{90.9}
\@namedef{cf@cf.fourseam.delta}{\ensuremath{-}4.5}
\@namedef{cf@cf.fourseam.delta.abs}{4.5}
\@namedef{cf@cf.fourseam.neutral}{95.5}
\@namedef{cf@cf.sinker.cdmx}{89.4}
\@namedef{cf@cf.sinker.delta}{\ensuremath{-}2.1}
\@namedef{cf@cf.sinker.delta.abs}{2.1}
\@namedef{cf@cf.sinker.neutral}{91.5}
\@namedef{cf@cf.slider.cdmx}{109.0}
\@namedef{cf@cf.slider.delta}{\ensuremath{-}0.3}
\@namedef{cf@cf.slider.delta.abs}{0.3}
\@namedef{cf@cf.slider.neutral}{109.3}
\@namedef{cf@cf.splitter.cdmx}{107.5}
\@namedef{cf@cf.splitter.delta}{\ensuremath{-}1.3}
\@namedef{cf@cf.splitter.delta.abs}{1.3}
\@namedef{cf@cf.splitter.neutral}{108.7}
\@namedef{cf@datos.juegos}{2,127}
\@namedef{cf@datos.pitcheos}{634,624}
\@namedef{cf@datos.pitchers}{1,135}
\@namedef{cf@diag.loc.excess}{0.13}
\@namedef{cf@diag.loc.r2oos}{0.31}
\@namedef{cf@diag.loc.r2ref}{0.18}
\@namedef{cf@escala.loc.sigma}{0.0742}
\@namedef{cf@escala.pit.sigma}{0.0774}
\@namedef{cf@escala.stuff.punto}{0.00054}
\@namedef{cf@escala.stuff.sigma}{0.0054}
\@namedef{cf@h1.changeup.hr}{+0.09}
\@namedef{cf@h1.changeup.hr.abs}{0.09}
\@namedef{cf@h1.changeup.hrp}{0.52}
\@namedef{cf@h1.changeup.n}{541}
\@namedef{cf@h1.changeup.tib}{\ensuremath{-}4.5}
\@namedef{cf@h1.changeup.tib.abs}{4.5}
\@namedef{cf@h1.changeup.tibabs}{4.5}
\@namedef{cf@h1.changeup.whiff}{\ensuremath{-}1.96}
\@namedef{cf@h1.changeup.whiff.abs}{1.96}
\@namedef{cf@h1.changeup.whiffp}{\ensuremath{<}0.01}
\@namedef{cf@h1.changeup.zs}{+0.89}
\@namedef{cf@h1.changeup.zs.abs}{0.89}
\@namedef{cf@h1.curveball.hr}{+0.38}
\@namedef{cf@h1.curveball.hr.abs}{0.38}
\@namedef{cf@h1.curveball.hrp}{0.28}
\@namedef{cf@h1.curveball.n}{363}
\@namedef{cf@h1.curveball.tib}{\ensuremath{-}3.7}
\@namedef{cf@h1.curveball.tib.abs}{3.7}
\@namedef{cf@h1.curveball.tibabs}{3.7}
\@namedef{cf@h1.curveball.whiff}{\ensuremath{-}0.27}
\@namedef{cf@h1.curveball.whiff.abs}{0.27}
\@namedef{cf@h1.curveball.whiffp}{0.75}
\@namedef{cf@h1.curveball.zs}{+1.05}
\@namedef{cf@h1.curveball.zs.abs}{1.05}
\@namedef{cf@h1.cutter.hr}{+0.18}
\@namedef{cf@h1.cutter.hr.abs}{0.18}
\@namedef{cf@h1.cutter.hrp}{0.62}
\@namedef{cf@h1.cutter.n}{92}
\@namedef{cf@h1.cutter.tib}{\ensuremath{-}2.7}
\@namedef{cf@h1.cutter.tib.abs}{2.7}
\@namedef{cf@h1.cutter.tibabs}{2.7}
\@namedef{cf@h1.cutter.whiff}{+0.88}
\@namedef{cf@h1.cutter.whiff.abs}{0.88}
\@namedef{cf@h1.cutter.whiffp}{0.48}
\@namedef{cf@h1.cutter.zs}{+1.08}
\@namedef{cf@h1.cutter.zs.abs}{1.08}
\@namedef{cf@h1.fourseam.hr}{+0.09}
\@namedef{cf@h1.fourseam.hr.abs}{0.09}
\@namedef{cf@h1.fourseam.hrp}{0.32}
\@namedef{cf@h1.fourseam.n}{642}
\@namedef{cf@h1.fourseam.tib}{\ensuremath{-}4.7}
\@namedef{cf@h1.fourseam.tib.abs}{4.7}
\@namedef{cf@h1.fourseam.tibabs}{4.7}
\@namedef{cf@h1.fourseam.whiff}{\ensuremath{-}1.85}
\@namedef{cf@h1.fourseam.whiff.abs}{1.85}
\@namedef{cf@h1.fourseam.whiffp}{\ensuremath{<}0.01}
\@namedef{cf@h1.fourseam.zs}{+1.35}
\@namedef{cf@h1.fourseam.zs.abs}{1.35}
\@namedef{cf@h1.sinker.hr}{+0.18}
\@namedef{cf@h1.sinker.hr.abs}{0.18}
\@namedef{cf@h1.sinker.hrp}{0.12}
\@namedef{cf@h1.sinker.n}{409}
\@namedef{cf@h1.sinker.tib}{\ensuremath{-}4.8}
\@namedef{cf@h1.sinker.tib.abs}{4.8}
\@namedef{cf@h1.sinker.tibabs}{4.8}
\@namedef{cf@h1.sinker.whiff}{\ensuremath{-}1.07}
\@namedef{cf@h1.sinker.whiff.abs}{1.07}
\@namedef{cf@h1.sinker.whiffp}{0.02}
\@namedef{cf@h1.sinker.zs}{+1.34}
\@namedef{cf@h1.sinker.zs.abs}{1.34}
\@namedef{cf@h1.slider.hr}{+0.16}
\@namedef{cf@h1.slider.hr.abs}{0.16}
\@namedef{cf@h1.slider.hrp}{0.22}
\@namedef{cf@h1.slider.n}{648}
\@namedef{cf@h1.slider.tib}{\ensuremath{-}2.0}
\@namedef{cf@h1.slider.tib.abs}{2.0}
\@namedef{cf@h1.slider.tibabs}{2.0}
\@namedef{cf@h1.slider.whiff}{\ensuremath{-}0.96}
\@namedef{cf@h1.slider.whiff.abs}{0.96}
\@namedef{cf@h1.slider.whiffp}{0.04}
\@namedef{cf@h1.slider.zs}{+0.93}
\@namedef{cf@h1.slider.zs.abs}{0.93}
\@namedef{cf@h1.splitter.hr}{\ensuremath{-}0.11}
\@namedef{cf@h1.splitter.hr.abs}{0.11}
\@namedef{cf@h1.splitter.hrp}{0.75}
\@namedef{cf@h1.splitter.ivbp}{0.18}
\@namedef{cf@h1.splitter.n}{35}
\@namedef{cf@h1.splitter.p}{0.002}
\@namedef{cf@h1.splitter.tib}{\ensuremath{-}3.2}
\@namedef{cf@h1.splitter.tib.abs}{3.2}
\@namedef{cf@h1.splitter.tibabs}{3.2}
\@namedef{cf@h1.splitter.whiff}{\ensuremath{-}2.04}
\@namedef{cf@h1.splitter.whiff.abs}{2.04}
\@namedef{cf@h1.splitter.whiffp}{0.41}
\@namedef{cf@h1.splitter.zs}{+0.51}
\@namedef{cf@h1.splitter.zs.abs}{0.51}
\@namedef{cf@h2.ci}{[0.65, 0.94]}
\@namedef{cf@h2.est}{0.79}
\@namedef{cf@h2.n}{754}
\@namedef{cf@h3.ci}{[\ensuremath{-}0.107, 0.100]}
\@namedef{cf@h3.est}{\ensuremath{-}0.010}
\@namedef{cf@h3.est.abs}{0.010}
\@namedef{cf@h3.hr.est}{\ensuremath{-}0.0015}
\@namedef{cf@h3.hr.est.abs}{0.0015}
\@namedef{cf@h3.hr.p}{0.90}
\@namedef{cf@h3.n}{792}
\@namedef{cf@h3.p}{0.83}
\@namedef{cf@h4.n}{892}
\@namedef{cf@h4.whiff.hb.p}{0.15}
\@namedef{cf@h4.whiff.hb.per100}{0.1809}
\@namedef{cf@h4.whiff.hb.prel}{=0.15}
\@namedef{cf@h4.whiff.ivb.p}{0.47}
\@namedef{cf@h4.whiff.ivb.per100}{\ensuremath{-}0.1859}
\@namedef{cf@h4.whiff.ivb.per100.abs}{0.1859}
\@namedef{cf@h4.whiff.ivb.prel}{=0.47}
\@namedef{cf@h4.whiff.s2b.p}{0.04}
\@namedef{cf@h4.whiff.s2b.per100}{\ensuremath{-}0.3603}
\@namedef{cf@h4.whiff.s2b.per100.abs}{0.3603}
\@namedef{cf@h4.whiff.s2b.prel}{=0.04}
\@namedef{cf@h4.whiff.spin.p}{0.03}
\@namedef{cf@h4.whiff.spin.per100}{0.0032}
\@namedef{cf@h4.whiff.spin.prel}{=0.03}
\@namedef{cf@h4.xrv.hb.p}{0.69}
\@namedef{cf@h4.xrv.hb.per100}{\ensuremath{-}0.0460}
\@namedef{cf@h4.xrv.hb.per100.abs}{0.0460}
\@namedef{cf@h4.xrv.hb.prel}{=0.69}
\@namedef{cf@h4.xrv.ivb.p}{0.31}
\@namedef{cf@h4.xrv.ivb.per100}{\ensuremath{-}0.0977}
\@namedef{cf@h4.xrv.ivb.per100.abs}{0.0977}
\@namedef{cf@h4.xrv.ivb.prel}{=0.31}
\@namedef{cf@h4.xrv.s2b.p}{0.74}
\@namedef{cf@h4.xrv.s2b.per100}{\ensuremath{-}0.0443}
\@namedef{cf@h4.xrv.s2b.per100.abs}{0.0443}
\@namedef{cf@h4.xrv.s2b.prel}{=0.74}
\@namedef{cf@h4.xrv.spin.p}{\ensuremath{<}0.01}
\@namedef{cf@h4.xrv.spin.per100}{0.0028}
\@namedef{cf@h4.xrv.spin.prel}{\ensuremath{<}0.01}
\@namedef{cf@lb.corr.fsuso}{\ensuremath{-}0.51}
\@namedef{cf@lb.corr.fsuso.abs}{0.51}
\@namedef{cf@lb.delta.max}{+0.47}
\@namedef{cf@lb.delta.max.abs}{0.47}
\@namedef{cf@lb.delta.mean}{\ensuremath{-}2.38}
\@namedef{cf@lb.delta.mean.abs}{2.38}
\@namedef{cf@lb.delta.median}{\ensuremath{-}2.14}
\@namedef{cf@lb.delta.median.abs}{2.14}
\@namedef{cf@lb.delta.min}{\ensuremath{-}8.60}
\@namedef{cf@lb.delta.min.abs}{8.60}
\@namedef{cf@lb.delta.pos}{10}
\@namedef{cf@lb.fsalto.delta}{\ensuremath{-}4.04}
\@namedef{cf@lb.fsalto.delta.abs}{4.04}
\@namedef{cf@lb.fsalto.n}{70}
\@namedef{cf@lb.fsalto.umbral}{60}
\@namedef{cf@lb.fsbajo.delta}{\ensuremath{-}1.62}
\@namedef{cf@lb.fsbajo.delta.abs}{1.62}
\@namedef{cf@lb.fsbajo.n}{172}
\@namedef{cf@lb.fsbajo.umbral}{30}
\@namedef{cf@lb.min}{500}
\@namedef{cf@lb.n}{464}
\@namedef{cf@relabel.ch.auto}{\ensuremath{-}3.44}
\@namedef{cf@relabel.ch.auto.abs}{3.44}
\@namedef{cf@relabel.ch.inv}{\ensuremath{-}1.62}
\@namedef{cf@relabel.ch.inv.abs}{1.62}
\@namedef{cf@relabel.fc.auto}{+2.15}
\@namedef{cf@relabel.fc.auto.abs}{2.15}
\@namedef{cf@relabel.fc.inv}{+0.12}
\@namedef{cf@relabel.fc.inv.abs}{0.12}
\@namedef{cf@relabel.si.auto}{\ensuremath{-}1.58}
\@namedef{cf@relabel.si.auto.abs}{1.58}
\@namedef{cf@relabel.si.inv}{\ensuremath{-}0.23}
\@namedef{cf@relabel.si.inv.abs}{0.23}
\@namedef{cf@relabel.sl.auto}{+3.10}
\@namedef{cf@relabel.sl.auto.abs}{3.10}
\@namedef{cf@relabel.sl.inv}{+1.75}
\@namedef{cf@relabel.sl.inv.abs}{1.75}
\@namedef{cf@rho.cdmx}{0.852}
\@namedef{cf@rho.ext.max}{0.867}
\@namedef{cf@rho.ext.min}{0.852}
\@namedef{cf@rho.ext.p90max}{0.919}
\@namedef{cf@rho.ext.pitcheos}{183,175}
\@namedef{cf@rho.med.max}{0.935}
\@namedef{cf@rho.med.min}{0.893}
\@namedef{cf@rho.med.p10min}{0.854}
\@namedef{cf@rho.med.pitcheos}{154,774}
\@namedef{cf@rho.no.max}{1.015}
\@namedef{cf@rho.no.min}{0.989}
\@namedef{cf@rho.no.pitcheos}{296,675}
\@namedef{cf@riesgo.hr.familias}{7}
\@namedef{cf@riesgo.hr.signif}{0}
\@namedef{cf@riesgo.hr.signif.txt}{ninguna diferencia es significativa}
\@namedef{cf@riesgo.hr.suben}{6}
\@namedef{cf@sc.fs.best}{95.7}
\@namedef{cf@sc.fs.gain}{+5.4}
\@namedef{cf@sc.fs.gain.abs}{5.4}
\@namedef{cf@sc.fs.nrows}{25}
\@namedef{cf@sc.fs.p25}{88.6}
\@namedef{cf@sc.fs.p75}{92.3}
\@namedef{cf@sc.peers}{73}
\@namedef{cf@sc.relh.tol}{0.35}
\@namedef{cf@sc.rep.changeup}{103.4}
\@namedef{cf@sc.rep.changeup.n}{30}
\@namedef{cf@sc.rep.curveball}{107.6}
\@namedef{cf@sc.rep.curveball.n}{9}
\@namedef{cf@sc.rep.cutter}{99.4}
\@namedef{cf@sc.rep.cutter.n}{8}
\@namedef{cf@sc.rep.sinker}{88.2}
\@namedef{cf@sc.rep.sinker.n}{29}
\@namedef{cf@sc.rep.slider}{111.3}
\@namedef{cf@sc.rep.slider.n}{49}
\@namedef{cf@sc.rep.splitter}{106.1}
\@namedef{cf@sc.rep.splitter.n}{4}
\@namedef{cf@sc.slot.tol}{6}
\@namedef{cf@sc.velo.tol}{2}
\@namedef{cf@stab.loc.r}{0.61}
\@namedef{cf@stab.loc.rho}{0.60}
\@namedef{cf@stab.n}{202}
\@namedef{cf@stab.pit.r}{0.60}
\@namedef{cf@stab.pit.rho}{0.61}
\@namedef{cf@stab.stuff.r}{0.74}
\@namedef{cf@stab.stuff.rho}{0.76}
\@namedef{cf@stuff.alpha}{0.01}
\@namedef{cf@stuff.lgbm}{num\_leaves=7, min\_child\_samples=1000, reg\_lambda=300.0}
\@namedef{cf@val.loc.gb.mejor}{84}
\@namedef{cf@val.loc.gb.peor}{96}
\@namedef{cf@val.loc.oop.max}{10.6}
\@namedef{cf@val.loc.oop.min}{9.3}
\@namedef{cf@val.park.base.r}{0.162}
\@namedef{cf@val.park.base.rmse}{0.2115}
\@namedef{cf@val.park.loc.r}{0.320}
\@namedef{cf@val.park.loc.rmse}{0.2030}
\@namedef{cf@val.park.pit.r}{0.322}
\@namedef{cf@val.park.pit.rmse}{0.2027}
\@namedef{cf@val.park.stuff.r}{0.164}
\@namedef{cf@val.park.stuff.rmse}{0.2115}
\@namedef{cf@val.pit.base.r}{0.157}
\@namedef{cf@val.pit.base.rmse}{0.2081}
\@namedef{cf@val.pit.ll.max}{39.7}
\@namedef{cf@val.pit.ll.min}{9.7}
\@namedef{cf@val.pit.loc.r}{0.328}
\@namedef{cf@val.pit.loc.rmse}{0.1991}
\@namedef{cf@val.pit.pit.r}{0.332}
\@namedef{cf@val.pit.pit.rmse}{0.1988}
\@namedef{cf@val.pit.stuff.r}{0.159}
\@namedef{cf@val.pit.stuff.rmse}{0.2081}
\@namedef{cf@val.sea.base.r}{0.141}
\@namedef{cf@val.sea.base.rmse}{0.2176}
\@namedef{cf@val.sea.loc.r}{0.329}
\@namedef{cf@val.sea.loc.rmse}{0.2075}
\@namedef{cf@val.sea.pit.r}{0.333}
\@namedef{cf@val.sea.pit.rmse}{0.2072}
\@namedef{cf@val.sea.stuff.r}{0.144}
\@namedef{cf@val.sea.stuff.rmse}{0.2175}
\@namedef{cf@val.stuff.cs.max}{0.0}
\@namedef{cf@val.stuff.cs.min}{\ensuremath{-}0.5}
\@namedef{cf@val.stuff.cs.min.abs}{0.5}
\@namedef{cf@val.stuff.ganados}{18}
\@namedef{cf@val.stuff.gb.max}{1.2}
\@namedef{cf@val.stuff.gb.min}{1.1}
\@namedef{cf@val.stuff.total}{21}
\@namedef{cf@val.stuff.whiff.max}{1.9}
\@namedef{cf@val.stuff.whiff.min}{1.7}
\@namedef{cf@vv.icc}{0.997}
\@namedef{cf@vv.mae.ann}{12.2}
\@namedef{cf@vv.mae.ann.n}{3}
\@namedef{cf@vv.mae.det}{6.4}
\@namedef{cf@vv.mae.det.n}{4}
\@namedef{cf@vv.n}{4}
\@namedef{cf@vv.rel.ms}{177}
\@namedef{cf@vv.window}{15}
\makeatother


\newcommand{\splus}{\text{Stuff+}}
\newcommand{\lplus}{\text{Location+}}
\newcommand{\pplus}{\text{Pitching+}}
\newcommand{\rhohat}{\hat{\rho}}
\newcommand{\tabla}[1]{\input{generated/#1.tex}}

\newcommand{\caja}[2]{%
  \par\smallskip\noindent\colorbox{grisfondo}{\parbox{\dimexpr\linewidth-2\fboxsep}{%
    \textbf{\color{diablos}#1}\\[2pt]#2}}\par\smallskip}

\pagestyle{fancy}
\fancyhf{}
\fancyhead[L]{\small\color{diablos}\textbf{Diablos AeroVision}}
\fancyhead[R]{\small ISAC Sports Analytics Hackathon 2026}
\fancyfoot[C]{\small\thepage}
\renewcommand{\headrulewidth}{0.4pt}


\title{\textbf{Diablos AeroVision}\\[4pt]
\large Perfil físico ideal para el Harp Helú, recomendaciones operativas\\ y tres casos para la entrevista}
\author{Entregables 2 y 3 · Reto 2 Diablos Rojos · ISAC Sports Analytics Hackathon 2026}
\date{22 de octubre de 2026}

\begin{document}
\maketitle
\thispagestyle{fancy}

\caja{Resumen para el cuerpo técnico}{%
\begin{itemize}[leftmargin=12pt]
  \item En el aire de la CDMX \textbf{todos los pitcheos se mueven menos}. Dentro del mismo pitcher, la recta de cuatro costuras pierde \cf{h1.fourseam.tibabs}~in de quiebre total y el sinker \cf{h1.sinker.tibabs}~in; el slider solo \cf{h1.slider.tibabs}~in.
  \item \textbf{El pitcheo que más valor pierde es la recta de cuatro costuras}: en promedio \cf{cf.fourseam.delta.abs} puntos de \splus\ al llevarla a la densidad de la CDMX. Slider (\cf{cf.slider.delta}), cutter (\cf{cf.cutter.delta}) y curva (\cf{cf.curveball.delta}) casi no cambian.
  \item Los pitchers que usan la recta en al menos el \cf{lb.fsalto.umbral}\% de sus lanzamientos pierden en promedio \cf{lb.fsalto.delta.abs} puntos; los que la usan menos del \cf{lb.fsbajo.umbral}\%, \cf{lb.fsbajo.delta.abs}.
  \item \textbf{La mayor parte del daño atribuible es del parque, no del pitcheo.} La bola bateada vuela más: es la misma calidad de contacto, que en altitud produce más daño. No hay que castigar a un pitcher por eso al evaluarlo.
  \item \textbf{Cómo usarlo:} evaluar con \splus\ CDMX y su $\Delta$, no con \splus\ neutral; en casa, mover el uso hacia los pitcheos resistentes; rediseñar con el Pitch Sculptor solo dentro de lo que pitchers con la misma mecánica ya lanzan.
\end{itemize}}

%=====================================================================
\section{Perfil físico y de arsenal ideal en el Harp Helú}

El perfil se construye solo con lo que los datos LMB confirman (ver el documento técnico):

\begin{enumerate}
  \item \textbf{El arma principal no es la recta de cuatro costuras.} Es la familia con la mayor pérdida de \splus\ y una de las de mayor pérdida de quiebre (el sinker pierde algo más). En el leaderboard (\cf{lb.n} temporadas de pitcher con al menos \cf{lb.min} pitcheos), el uso de la recta de cuatro costuras correlaciona $r=\cf{lb.corr.fsuso}$ con el $\Delta\splus_{\text{CDMX}}$.
  \item \textbf{Arsenal cargado a slider, cutter y curva.} Son las familias que conservan su \splus\ (Tabla~\ref{tab:cf}). El slider, además, es la familia que menos quiebre total pierde (\cf{h1.slider.tibabs}~in).
  \item \textbf{Spin alto.} En \emph{Extreme Altitude}, de las cuatro variables de la regresión conjunta (spin, eficiencia de giro, IVB y HB), el \texttt{SpinRate} bruto es la única asociada con más valor por pitcheo ($p\cf{h4.xrv.spin.prel}$); para whiff también lo es ($p\cf{h4.whiff.spin.prel}$), junto con \texttt{SpinToBreakRatio} ($p\cf{h4.whiff.s2b.prel}$, con signo negativo). Es una asociación entre pitchers dentro de altitud, no un efecto causal ni una comparación con el nivel del mar.
  \item \textbf{La velocidad no se pierde, se gana.} Con menos arrastre, la recta llega \cf{h1.fourseam.zs}~mph y el sinker \cf{h1.sinker.zs}~mph más rápido al plato. La velocidad es un activo en la CDMX, pero no compensa la pérdida de movimiento en el \splus.
  \item \textbf{El cambio y el splitter no son la apuesta segura.} El modelo les resta \cf{h2.est} puntos menos que a la recta, la curva y el slider del mismo pitcher (H2), pero pierden más \splus\ que el cutter y la curva (\cf{cf.changeup.delta} y \cf{cf.splitter.delta} frente a \cf{cf.cutter.delta} y \cf{cf.curveball.delta}). En lo observado, el cambio pierde \cf{h1.changeup.tibabs}~in de quiebre, casi como la recta, y \cf{h1.changeup.whiff.abs}~pp de whiff, la mayor caída significativa.
  \item \textbf{Ser pitcher de rodados no es, por sí solo, una ventaja demostrada.} La tasa basal de rodados no predice quién sufre menos en altitud ($p=\cf{h3.p}$, H3).
\end{enumerate}

\begin{table}[H]\centering\small
\caption{Qué gana y qué pierde \splus\ en la CDMX por familia (contrafactual del modelo; media ponderada por pitcheos en toda la LMB).}\label{tab:cf}
\tabla{tab_contrafactual}
\end{table}

%=====================================================================
\section{Perfiles que ganan y perfiles que pierden}

\paragraph{Ganan valor relativo:} los arsenales con poca dependencia de la recta de cuatro costuras, donde el slider, el cutter y la curva pesan más (Tabla~\ref{tab:cf}). De las \cf{lb.n} temporadas analizadas, solo \cf{lb.delta.pos} tienen $\Delta\splus_{\text{CDMX}}$ positivo (máximo \cf{lb.delta.max}). ``Ganar'' en la CDMX significa, en la práctica, \emph{no perder}, mientras la mediana de la liga pierde \cf{lb.delta.median.abs} puntos.

\paragraph{Pierden valor:} la recta de cuatro costuras y el sinker, las dos familias con mayor pérdida de quiebre y de \splus. El caso extremo del leaderboard pierde \cf{lb.delta.min.abs} puntos.

\paragraph{Riesgo de contacto duro y HR.} Dentro del mismo pitcher, la tasa de HR por pitcheo sube en altitud en \cf{riesgo.hr.suben} de las \cf{riesgo.hr.familias} familias, con \cf{riesgo.hr.signif} diferencias significativas al 5\% por separado (Tabla~\ref{tab:riesgo}). El daño por contacto sí aparece en conjunto: el canal parque resta \cf{canal.park.per100.abs} carreras por cada 100 pitcheos (IC por pitcheo \cf{canal.park.ci}) a igual velocidad y ángulo de salida. En contraste, la pérdida de whiff sí es clara en la recta (\cf{h1.fourseam.whiff}~pp), el cambio (\cf{h1.changeup.whiff}~pp), el sinker y el slider.

\begin{table}[H]\centering\small
\caption{Riesgo por familia en \emph{Extreme} frente a \emph{No Altitude} (dentro del mismo pitcher, en puntos porcentuales por pitcheo).}\label{tab:riesgo}
\tabla{tab_riesgo}
\end{table}

%=====================================================================
\section{Cómo perfilar a un prospecto o agente libre}

\begin{enumerate}
  \item \textbf{Con TrackMan:} cargar sus pitcheos en el simulador. Comparar el \splus\ neutral con el CDMX por pitcheo y revisar si su arsenal depende de la recta.
  \item \textbf{Sin TrackMan (VeloVision, experimental):} grabar un bullpen de lado, marcar la escala, \emph{confirmar el frame de soltura} y obtener el slot del brazo. La herramienta bloquea el resultado si el video no cumple la puerta de calidad y devuelve rangos de \splus\ CDMX por familia de pitcheo, como referencia de lo que logran pitchers de la LMB con esa mecánica, no como el nivel del prospecto. En el piloto con 2 anotadores ($n=\cf{vv.n}$) el slot por video difirió de la referencia humana ${\sim}\cf{vv.mae.det}$--$\cf{vv.mae.ann}^\circ$ (MAE); con $n$ tan pequeño es un orden de magnitud, y el error frente a TrackMan no está medido.
  \item \textbf{Rediseño (Pitch Sculptor):} buscar, entre pitchers reales con el mismo slot ($\pm\cf{sc.slot.tol}^\circ$), altura de soltura ($\pm\cf{sc.relh.tol}$~ft) y velocidad de recta ($\pm\cf{sc.velo.tol}$~mph), qué pitcheos y formas rinden más en la CDMX. Las diferencias son observadas frente a pares, no ganancias causales garantizadas.
  \item \textbf{Decisión:} priorizar \pplus\ CDMX para pronosticar rendimiento y \splus\ CDMX junto con su $\Delta$ para entender \emph{por qué}. Cuando haya varias temporadas, revisar que el $\Delta$ se repita (no lo hemos medido de forma sistemática).
\end{enumerate}

\paragraph{Staff de Diablos.} Los datos no traen equipo ni estadio y los IDs son anónimos, así que no podemos identificar a los pitchers de Diablos. Las recomendaciones se expresan por perfil y aplican a cualquier pitcher que el club cargue en el simulador.

%=====================================================================
\section{Tres casos para la entrevista de 10 minutos}

Los tres pitchers son reales (IDs anónimos del dataset). Los lugares ``de \cf{lb.n}'' ordenan pitcher-temporadas de 2024 a 2026 con al menos \cf{lb.min} pitcheos, no una liga de un solo año. Los scores de 2024--2025 son dentro de muestra (esas temporadas se usaron para entrenar); solo los de 2026 son fuera de muestra.

%---------------------------------------------------------------------
\subsection*{Caso A · Fichaje: el agente libre que las métricas neutrales subestiman}

\textbf{Situación.} \cf{ca2.id}, \cf{ca2.mano}, \cf{ca2.velo}~mph de velocidad media, arsenal de slider (\cf{ca.slider.uso}\%) y cambio (\cf{ca.changeup.uso}\%). Su \splus\ neutral es modesto (\cf{ca2.sn} en \cf{ca2.year}; lugar \cf{ca2.rank.sn} de \cf{lb.n}). Un scout de nivel del mar lo descartaría.

\textbf{Qué dice el modelo.}
\begin{itemize}
  \item No pierde nada en la CDMX: $\Delta\splus = \cf{ca1.sd}$ en \cf{ca1.year} y \cf{ca2.sd} en \cf{ca2.year}. Está entre los más resistentes de la liga (lugar \cf{ca1.rank.delta} y \cf{ca2.rank.delta} de \cf{lb.n}) dos años seguidos.
  \item Su valor está en el comando: \lplus\ \cf{ca1.ln} y \cf{ca2.ln}; \pplus\ CDMX \cf{ca1.pc} y \cf{ca2.pc}, ambos por encima del promedio de la liga.
  \item El slider pasa de \cf{ca.slider.sn} a \cf{ca.slider.sc} en la CDMX y el cambio de \cf{ca.changeup.sn} a \cf{ca.changeup.sc} (Tabla~\ref{tab:ca}). Además, \cf{ca2.gb}\% de sus pelotas en juego en \cf{ca2.year} fueron rodados, y su valor observado fue mejor en \emph{Extreme} que en \emph{No Altitude} (\cf{ca1.ct} y \cf{ca2.ct} carreras por 100 pitcheos).
\end{itemize}

\begin{table}[H]\centering\small
\caption{Caso A: arsenal \cf{ca2.year}.}\label{tab:ca}
\tabla{tab_ca_arsenal}
\end{table}

\textbf{Decisión.} Fichar: es un perfil hecho para el Harp Helú, y un análisis con métricas neutrales lo subestima. \textbf{Riesgo:} su \splus\ absoluto es bajo; si pierde comando, no tiene la física para compensarlo.

%---------------------------------------------------------------------
\subsection*{Caso B · Alerta: la trampa del nivel del mar}

\textbf{Situación.} \cf{cb.id}, \cf{cb.mano}, \cf{cb.velo}~mph. En \cf{cb.year} su \splus\ neutral es \cf{cb.sn}: lugar \cf{cb.rank.sn} de \cf{lb.n} (percentil \cf{cb.pctl}). En papel, un brazo élite.

\textbf{Qué dice el modelo.}
\begin{itemize}
  \item En la CDMX baja a \cf{cb.sc} ($\Delta=\cf{cb.sd}$), el lugar \cf{cb.rank.delta} de \cf{lb.n} en resistencia a la altitud. \pplus\ pasa de \cf{cb.pn} a \cf{cb.pc}.
  \item La causa es una: lanza la recta de cuatro costuras el \cf{cb.fourseam.uso}\% de las veces, y esa recta cae de \cf{cb.fourseam.sn} a \cf{cb.fourseam.sc}. Su slider casi no cambia (\cf{cb.slider.sn} $\to$ \cf{cb.slider.sc}).
  \item \textbf{Lo que el modelo no dice:} en su propia muestra de \cf{cb.year}, su valor observado fue \emph{mejor} en \emph{Extreme} que en \emph{No Altitude} (\cf{cb.ct} carreras por 100 pitcheos). \splus\ describe la física de sus pitcheos, no ese resultado, que también depende de ubicación, rivales y azar. La alerta es sobre el riesgo físico de su arsenal, no una predicción de que vaya a fracasar.
\end{itemize}

\begin{table}[H]\centering\small
\caption{Caso B: arsenal \cf{cb.year}.}\label{tab:cb}
\tabla{tab_cb_arsenal}
\end{table}

\textbf{Decisión.} No pagar por su \splus\ neutral sin revisar su rendimiento observado en altitud. Si se contrata, que sea con un plan explícito: en casa, invertir la mezcla hacia el slider (\cf{cb.slider.uso}\% de uso actual). \textbf{Riesgo:} un slider usado mucho más podría perder efectividad; habría que vigilarlo en el simulador y en el juego.

%---------------------------------------------------------------------
\subsection*{Caso C · Rediseño con Pitch Sculptor y plan casa contra gira}

\textbf{Situación.} \cf{cc.id}, \cf{cc.mano}, slot de \cf{cc.slot}$^\circ$, \cf{cc.velo}~mph. Su recta (\cf{cc.fourseam.uso}\% de uso) es promedio a nivel del mar (\cf{cc.fourseam.sn}) pero cae a \cf{cc.fourseam.sc} en la CDMX ($\Delta=\cf{cc.fourseam.sd}$). Su curva (\cf{cc.curveball.sn} $\to$ \cf{cc.curveball.sc}) y su cambio (\cf{cc.changeup.sn} $\to$ \cf{cc.changeup.sc}) son sus mejores pitcheos en ambos entornos. Dentro de su propia muestra, su valor observado en \emph{Extreme} es \cf{cc.canaltotal.abs} carreras por 100 pitcheos menor que en \emph{No Altitude}.

\begin{table}[H]\centering\small
\caption{Caso C: arsenal \cf{cc.year}.}\label{tab:cc}
\tabla{tab_cc_arsenal}
\end{table}

\textbf{Qué dice el Pitch Sculptor.} Excluyéndolo a él del catálogo, encuentra \cf{sc.peers} pitchers reales con su misma mecánica. Las \cf{sc.fs.nrows} rectas (pitcher-temporada) más cercanas a la suya tienen un rango intercuartílico de \cf{sc.fs.p25}--\cf{sc.fs.p75} de \splus\ CDMX, y la mejor referencia factible llega a \cf{sc.fs.best} (\cf{sc.fs.gain} frente a la suya). La señal más fuerte está en el repertorio: \cf{sc.rep.slider.n} de esos pares lanzan slider, con mediana de \splus\ CDMX de \cf{sc.rep.slider}, y él no tiene slider (Tabla~\ref{tab:sc}).

\begin{table}[H]\centering\small
\caption{Caso C: qué lanzan los pares factibles y cuánto rinde en la CDMX.}\label{tab:sc}
\tabla{tab_sculptor}
\end{table}

\textbf{Decisión.}
\begin{itemize}
  \item \textbf{En casa (Harp Helú):} bajar el uso de la recta y subir la curva, que no pierde valor ($\Delta=\cf{cc.curveball.sd}$), y el cambio, que pierde poco frente a la recta (\cf{cc.changeup.sd} contra \cf{cc.fourseam.sd}). La recta queda como pitcheo de establecimiento, no como arma de dos strikes.
  \item \textbf{De gira (nivel del mar):} mantener la mezcla actual; ahí su recta rinde en el promedio de la liga.
  \item \textbf{Proyecto de desarrollo:} incorporar un slider, el pitcheo con mejor \splus\ CDMX entre pitchers con su misma mecánica.
\end{itemize}
\textbf{Riesgo:} el Sculptor reporta diferencias observadas entre pares reales, no ganancias causales. Aprender un pitcheo nuevo puede no reproducir la forma de los pares.

\end{document}
